\documentclass[11pt]{article}
\usepackage{amsmath,amssymb,amsthm}
\newtheorem{theorem}{Theorem}
\newtheorem{corollary}{Corollary}
\title{Moment valuations in the unperturbed Chebyshev tree experiment}
\author{}
\date{}
\begin{document}
\maketitle

\paragraph{Status and scope.}
The moment theorem and its consequences below have complete written
proofs. They have not been formalized in Lean. An accompanying exact
cyclotomic computation checks the finite cases $H=4,8,16$.
The result concerns one candidate mechanism, not all admissible heads.
In particular it does not establish either accepted global target.
It explains why generic fixed-numerator experiments with an unperturbed
composed pole set have much less branch diversity than the number $H!$
of permutations suggests. Numerators depending on the coalescence
parameter, moving query times, and perturbed poles are outside the
stated conclusion.

Let $H=2^d$, and let $P_H$ be the monic multiple of the Chebyshev
polynomial $T_H$. Write its distinct roots as
$a_i=\cos((2i+1)\pi/(2H))$, $0\leq i<H$, and put
$w_i=1/P_H'(a_i)$. For a subset $S$ of the roots define
\[
 \mu_k(S)=\sum_{i\in S}w_i a_i^k.
\]
For nonempty $S$, its valuation is the smallest $k\geq0$ for which
$\mu_k(S)\ne0$.

\begin{theorem}
Every nonempty subset has valuation $2^j-1$ for some $0\leq j\leq d$.
For $j<d$, the number of subsets with this valuation is
\[
 2^{H/2^j}-2^{H/2^{j+1}}.
\]
Only the full subset has valuation $H-1$.
\end{theorem}

\begin{proof}
First suppose $H\geq2$. Since $P_H$ is even, opposite roots have
opposite barycentric weights. The positive-root weights are linearly
independent over $\mathbb Q$. Indeed, up to a common nonzero constant
and individual signs, these weights are
$\sin((2i+1)\pi/(2H))$, $0\leq i<H/2$.
Let $\zeta=\exp(\pi\mathrm i/(2H))$. A rational dependence among these
sines would give a polynomial of degree less than $2H$ vanishing at
$\zeta$, with its nonzero coefficients in the distinct exponent positions
$2i+1$ and $2H-(2i+1)$. Such a polynomial is zero because the minimal
polynomial of $\zeta$ is $X^{2H}+1$. The latter irreducibility is the
usual power-of-two cyclotomic identity; alternatively, substituting
$X+1$ gives Eisenstein's criterion at $2$.
Consequently $\mu_0(S)=0$ if and only if $S$ is a union of opposite pairs.

For such a union, write $P_H(X)=Q(X^2)$ and replace every selected pair
$\{a,-a\}$ by $b=a^2$. Its two contributions give
\[
 \mu_{2l}(S)=0,\qquad
 \mu_{2l+1}(S)=\sum_{b\in S'}\frac{b^l}{Q'(b)}.
\]
The identity $T_H(X)=T_{H/2}(2X^2-1)$ shows that the roots of $Q$ are
the affine images $(1+u)/2$ of the roots of $P_{H/2}$. An affine change
of variable preserves the first nonzero moment index, since the
change between powers is triangular with nonzero diagonal. Thus if
$S'$ has valuation $v$ in the smaller Chebyshev problem, then $S$ has
valuation $2v+1$.

Induction, starting with $H=1$, proves that the possible indices are
$2^j-1$. More precisely, subsets whose moments vanish through index
$2^j-2$ are exactly unions of the $H/2^j$ blocks obtained after $j$
iterations of opposite-pair identification. There are
$2^{H/2^j}-1$ nonempty such unions. Subtracting the count at the next
level gives the stated formula. At the last level there is just the
full subset. Its first nonzero moment is $\mu_{H-1}=1$, by the leading
coefficient in Lagrange interpolation of $X^{H-1}$.
\end{proof}

\begin{corollary}
Fix a real polynomial $N$ of degree at most $H-1$, and finitely many
positive exterior query times. For $\varepsilon>0$ put
\[
 p_i=\varepsilon a_i,\qquad
 r_i=\frac{N(p_i)}{\prod_{j\ne i}(p_i-p_j)},\qquad
 F_S(t)=\sum_{i\in S}\frac{r_i}{t-p_i}.
\]
Let $N_{\leq k}$ be the Taylor polynomial of $N$ at zero through degree
$k$. Assume $N_{\leq 2^j-1}(t)\ne0$ for every query time and every
$0\leq j\leq d$. For every nonempty subset $S$ of valuation $k$,
\[
 F_S(t)=\varepsilon^{k+1-H}
 \mu_k(S)\frac{N_{\leq k}(t)}{t^{k+1}}
 +O(\varepsilon^{k+2-H}).
\]
For all sufficiently small positive $\varepsilon$, uniformly over every
subset and query, the displayed nonzero leading coefficient determines
the sign. Hence each fixed mask, as the pool is permuted, has at most
$2(d+1)$ possible sign profiles on all its queries.
\end{corollary}

\begin{proof}
Write $N(X)=\sum_l n_lX^l$ and expand
\[
 \frac{N(\varepsilon a_i)}{t-\varepsilon a_i}
 =\sum_{k\geq0}\varepsilon^k a_i^k
    \sum_{l=0}^{k}\frac{n_l}{t^{k-l+1}}.
\]
Here $n_l=0$ past the degree of $N$. The barycentric denominator is
$\varepsilon^{H-1}/w_i$. All lower moments vanish, and the coefficient
at the first remaining index is the stated Taylor truncation divided
by $t^{k+1}$. There are finitely many subsets and queries, so their
individual small-parameter sign bounds have a common positive minimum.
The previous theorem supplies only $d+1$ indices, with two possible
signs of $\mu_k(S)$.
\end{proof}

\paragraph{What this says about entropy.}
For a binary tree with $H-1$ proper right-child masks, a fixed numerator
therefore supplies at most $[2(d+1)]^{H-1}$ descendant patterns in this
generic germ regime. The bound is independent of how many fixed times
are queried at each mask. In fact the fixed-numerator bound improves to
$\exp(O(H))$: a mask of cardinality $M$ can only have valuation
$2^j-1$ when $2^j\leq M$, since it is then a union of blocks of size
$2^j$. There are $H/(2M)$ right-child masks of size $M$ in the complete
tree. Consequently the logarithm of the profile-count product is at most
\[
 \sum_{j=0}^{d-1}\frac{H}{2^{j+1}}\log(2(j+1))=O(H),
\]
because the corresponding infinite series converges. This refinement
was identified independently by the companion decoder audit.
Even allowing $N$ to vary, the full collection
of master signs $\operatorname{sgn}N_{\leq 2^j-1}(t)$ is determined by
at most $q(d+1)$ linear forms in the $H$ coefficients of $N$, where $q$
is the total number of queries. The elementary central-hyperplane
arrangement bound, followed by the mask choices, gives at most
\[
 2\left(\sum_{l=0}^{H-1}\binom{q(d+1)-1}{l}\right)
 [2(d+1)]^{H-1}
\]
patterns. For $q=O(Hd)$ this is $\exp(O(H\log(d+1)))$.
This is an upper bound for the described generic germs, not for all
finite values of $\varepsilon$ or all adaptive constructions.

There is a sharper bound that also excludes a more efficient query
selection within this same generic germ family.

\begin{corollary}[Linear VC dimension for the generic tree family]
Allow the numerator $N$ above and the pole permutation to vary. Fix a
complete binary tree of $H$ leaves. Queries may be made at arbitrary
fixed positive times, at the full mask or any proper right-child mask.
Retain the generic nonvanishing hypotheses on the relevant Taylor
truncations and take the limiting signs for sufficiently small
$\varepsilon$. The resulting class has VC dimension less than $8H$.
\end{corollary}
\begin{proof}
For each proper mask $S$, record only its first nonzero moment index
and the sign of that moment. A size-$2^j$ mask has at most $2(j+1)$
such states. The full mask always has index $H-1$ and moment $1$.
The number $K$ of state assignments over the proper masks is at most
\[
 K\leq\prod_{j=0}^{d-1}[2(j+1)]^{H/2^{j+1}}
 \leq 2^{2H},
\]
where the last inequality uses
$\log_2(2(j+1))\leq j+1$ and
$\sum_{j\geq0}(j+1)/2^{j+1}=2$.
Some assignments may not arise from any permutation; counting all
of them only increases the bound.

Conditional on one state assignment, every queried limiting sign is
the sign of a fixed linear form in the $H$ coefficients of $N$.
Indeed the positive factor $t^{-k-1}$ can be removed, and the remaining
form is the assigned moment sign times $N_{\leq k}(t)$.
On any $q\geq H$ queries, one such family has at most
$\sum_{j=0}^{H}\binom qj\leq(eq/H)^H$ strict sign patterns, by the
hyperplane arrangement bound. The union over all states consequently
has at most
\[
 4^H(eq/H)^H
\]
patterns. At $q=8H$ this is $(32e)^H<2^{8H}$, since $e<3$.
No set of $8H$ queries can therefore be shattered.
\end{proof}

The new corollary permits the coefficients of $N$ to depend on the
desired labels, but not to vary with $\varepsilon$ in a way that changes
the stated leading-term regime. Perturbed poles, query times moving
with $\varepsilon$, and vanishing Taylor truncations remain outside its
scope. It also does not bound the finite-$\varepsilon$ family uniformly.

\begin{corollary}[Arbitrary fixed scaled pole sets]
The same VC bound holds if the Chebyshev roots are replaced by any
$H$ distinct fixed real numbers $a_i$, with barycentric weights
$w_i=1/\prod_{j\ne i}(a_i-a_j)$, provided all relevant truncations
$N_{\leq k}(t)$, $0\leq k<H$, are nonzero at the fixed positive
query times. The poles remain $p_i=\varepsilon a_i$.
\end{corollary}
\begin{proof}
For a nonempty subset $S$ of size $M$, its first nonzero moment
$\sum_{i\in S}w_i a_i^k$ occurs for some $0\leq k<M$.
Otherwise the invertible $M\times M$ Vandermonde matrix on the
distinct points $(a_i)_{i\in S}$ would annihilate the nonzero vector
$(w_i)_{i\in S}$. Thus each size-$M$ mask has at most $2M$ possible
index/sign states. The full set has first index $H-1$ and moment $1$
by Lagrange interpolation, just as before.

The moment expansion used above did not require Chebyshev roots, so
conditional on the state assignment all query signs are still linear
in the coefficients of $N$. Finally, for $H=2^d$ the number of states
over the proper right-child masks is bounded by
\[
 \prod_{j=0}^{d-1}(2^{j+1})^{H/2^{j+1}}
 =2^{H\sum_{j=0}^{d-1}(j+1)/2^{j+1}}<4^H.
\]
The previous arrangement argument applies unchanged.
\end{proof}

In particular, using a different fixed real-rooted polynomial or
more polynomial-composition divisors does not by itself evade this
generic scaled-pole obstruction. A hierarchical pole deformation
such as $p_i=\varepsilon a_i+\varepsilon^2 b_i$ is outside it.

\paragraph{Reproducible computations.}
\texttt{tree-chebyshev-moment-check.py} works in
$\mathbb Q[\zeta]/(\zeta^{2H}+1)$ using integer coefficient vectors,
so its subset valuation counts are exact. It produces respectively
\[
\begin{array}{c|rrrrr}
H&k=0&k=1&k=3&k=7&k=15\\\hline
4&12&2&1&0&0\\
8&240&12&2&1&0\\
16&65280&240&12&2&1
\end{array}
\]
with one additional empty subset in each case.
The separate script \texttt{tree-cauchy-sixteen-entropy-probe.py} is
explicitly a high-precision numerical probe, not an exact certificate.
With five fixed-numerator pools and $500000$ sampled permutations per
pool, it found between $23521$ and $25635$ distinct descendant patterns
on $32$ queries. That observation motivated the moment audit; it is
not used in any proof above.
\end{document}
